\documentclass[11pt]{article}

% ── Encoding ──────────────────────────────────────────────────────────────────
\usepackage[T1]{fontenc}
\usepackage[utf8]{inputenc}

% ── Page layout ───────────────────────────────────────────────────────────────
\usepackage[margin=1in, headheight=14pt]{geometry}
\usepackage{fancyhdr}
\pagestyle{fancy}
\fancyhf{}
\fancyhead[L]{\small\textit{3psLCCA LaTeX Report Generator — Intern Screening Task}}
\fancyhead[R]{\small\thepage}
\fancyfoot[C]{\small\textit{Prepared by Shreyas Mene}}
\renewcommand{\headrulewidth}{0.4pt}
\renewcommand{\footrulewidth}{0.4pt}

% ── Colour palette ────────────────────────────────────────────────────────────
\usepackage{xcolor}
\definecolor{navyblue}{RGB}{20, 50, 100}
\definecolor{rowgray}{RGB}{245, 246, 248}
\definecolor{rowwhite}{RGB}{255, 255, 255}
\definecolor{codebg}{RGB}{248, 248, 248}

% ── Tables ────────────────────────────────────────────────────────────────────
\usepackage{booktabs}
\usepackage{array}
\usepackage{tabularx}
\usepackage{longtable}
\usepackage{colortbl}

% ── Math ──────────────────────────────────────────────────────────────────────
\usepackage{amsmath}
\usepackage{amssymb}
\usepackage{mathtools}

% ── Lists ─────────────────────────────────────────────────────────────────────
\usepackage{enumitem}
\setlist[itemize]{leftmargin=1.8em, itemsep=0.25em, topsep=0.4em}
\setlist[enumerate]{leftmargin=1.8em, itemsep=0.25em, topsep=0.4em}

% ── Figures & captions ────────────────────────────────────────────────────────
\usepackage{graphicx}
\usepackage{float}
\usepackage{caption}
\captionsetup{font=small, labelfont={bf,color=navyblue}, labelsep=period, skip=5pt}

% ── Section headings ──────────────────────────────────────────────────────────
\usepackage{titlesec}
\titleformat{\section}
  {\large\bfseries\color{navyblue}}{\thesection}{0.8em}{}
  [\vspace{-0.3em}\textcolor{navyblue}{\rule{\linewidth}{0.5pt}}]
\titleformat{\subsection}
  {\normalsize\bfseries\color{navyblue!80!black}}{\thesubsection}{0.7em}{}
\titlespacing*{\section}{0pt}{1.4em}{0.6em}
\titlespacing*{\subsection}{0pt}{1.0em}{0.4em}

% ── Hyperlinks ────────────────────────────────────────────────────────────────
\usepackage{hyperref}
\hypersetup{
  colorlinks=true, linkcolor=navyblue,
  urlcolor=navyblue, citecolor=navyblue,
  pdftitle={Intern Screening Task Report — 3psLCCA},
  pdfauthor={Shreyas Mene}
}

% ── Code listings ─────────────────────────────────────────────────────────────
\usepackage{listings}
\lstset{
  basicstyle=\small\ttfamily,
  backgroundcolor=\color{codebg},
  frame=single,
  framesep=4pt,
  rulecolor=\color{gray!35},
  breaklines=true,
  breakatwhitespace=true,
  xleftmargin=8pt,
  xrightmargin=4pt,
}

% ── Spacing ───────────────────────────────────────────────────────────────────
\setlength{\parindent}{0pt}
\setlength{\parskip}{0.6em}

% ── Column shorthands ─────────────────────────────────────────────────────────
\newcolumntype{L}[1]{>{\raggedright\arraybackslash}p{#1}}
\newcolumntype{C}[1]{>{\centering\arraybackslash}p{#1}}
% header-row helper: colour text white+bold for header cells
\newcommand{\hc}[1]{\color{white}\bfseries #1}

% ── Screenshot placeholder ────────────────────────────────────────────────────
\newcommand{\screenshotplaceholder}[2]{%
\begin{figure}[H]
  \centering
  \fcolorbox{navyblue!40}{rowgray}{%
    \begin{minipage}[c][0.26\textheight][c]{0.88\linewidth}
      \centering
      {\color{navyblue}\bfseries\large #1}\\[0.6em]
      {\color{gray!80}\itshape #2}
    \end{minipage}}
  \caption{#1}
\end{figure}}

% =============================================================================
\begin{document}

% ── Title page ────────────────────────────────────────────────────────────────
\begin{titlepage}
  \pagecolor{navyblue}
  \color{white}
  \centering
  \vspace*{3cm}

  {\fontsize{26}{32}\bfseries\selectfont Intern Screening Task Report\par}
  \vspace{0.5em}
  {\fontsize{17}{24}\bfseries\selectfont 3psLCCA \LaTeX{} Report Generator\par}

  \vspace{0.9em}
  \textcolor{white!50}{\rule{0.5\linewidth}{0.5pt}}
  \vspace{0.9em}

  {\large Life Cycle Cost Analysis (LCCA) Engine\\[0.3em]
   Implementation Summary\par}

  \vfill

  \renewcommand{\arraystretch}{1.5}
  \begin{tabular}{@{} L{3.5cm} L{7cm} @{}}
    \textbf{Project:}      & 3psLCCA Core Task \\
    \textbf{Repository:}   & \texttt{3psLCCA-core-task} \\
    \textbf{Prepared by:}  & Shreyas Mene \\
    \textbf{Date:}         & \today \\
  \end{tabular}
  \renewcommand{\arraystretch}{1}

  \vspace{2.5cm}
\end{titlepage}
\nopagecolor

% ── Table of Contents ─────────────────────────────────────────────────────────
\tableofcontents
\clearpage

% =============================================================================
\section{Task Context (From README)}

\subsection{Objective}
The task objective is to extend the existing 3psLCCA engine so it can generate a
\textbf{comprehensive, human-readable \LaTeX{} report} that explains lifecycle cost
calculations. The report is expected to:

\begin{itemize}
  \item Document every formula used.
  \item Explain each calculation in plain English.
  \item Show input values.
  \item Show substitution and step-wise derivation.
  \item Present final computed outputs in structured form.
\end{itemize}

\subsection{Background}
The project analyses bridge lifecycle costs over four stages:
\textbf{Initial}, \textbf{Use}, \textbf{Reconstruction}, and \textbf{End-of-Life}.
The core entry point is:

\begin{lstlisting}[language=Python]
# src/three_ps_lcca_core/core/main.py
run_full_lcc_analysis(...)
\end{lstlisting}

\screenshotplaceholder
  {Task Statement Reference (Add Screenshot)}
  {Insert screenshot of README objective and report-structure requirements.}

% =============================================================================
\section{Required Functional Change}

\subsection{Function Signature Requirement}
As required by the task, the analysis entry point supports:

\begin{lstlisting}[language=Python]
def run_full_lcc_analysis(
    input_data,
    construction_costs,
    debug=False,
    latex_report=False,
    latex_output_path=None,
):
\end{lstlisting}

\subsection{Behaviour Contract}

\renewcommand{\arraystretch}{1.45}
\begin{table}[H]
  \centering
  \caption{Parameter behaviour contract for \texttt{run\_full\_lcc\_analysis()}.}
  \small
  \begin{tabularx}{\linewidth}{%
      >{\columncolor{rowgray}}L{0.28\linewidth}%
      >{\columncolor{rowwhite}}L{0.22\linewidth}%
      >{\columncolor{rowgray}}X}
    \toprule
    \rowcolor{navyblue}
    \hc{Parameter} & \hc{Value} & \hc{Expected Behaviour} \\
    \midrule
    \texttt{latex\_report}      & \texttt{False}       & Existing behaviour remains unchanged. \\[0.15em]
    \texttt{latex\_report}      & \texttt{True}        & Generates a \texttt{.tex} report file. \\[0.15em]
    \texttt{latex\_output\_path}& \texttt{None}        & Saves output as \texttt{LCCA\_Report.tex}. \\[0.15em]
    \texttt{latex\_output\_path}& \textit{custom path} & Saves report at the provided path. \\
    \bottomrule
  \end{tabularx}
\end{table}
\renewcommand{\arraystretch}{1}

\textbf{Important rule:} When \texttt{latex\_report=True}, debug data must be internally
available for report creation, but debug JSON file dumping should occur only if the
caller explicitly sets \texttt{debug=True}.

\screenshotplaceholder
  {Main Entry-Point Update (Add Screenshot)}
  {Insert screenshot of \texttt{run\_full\_lcc\_analysis()} showing debug/latex behaviour handling.}

% =============================================================================
\section{Code Architecture and Logic Flow}

\subsection{High-Level Pipeline}
\begin{enumerate}
  \item Normalise and validate input (dictionary or \texttt{InputGlobalMetaData}).
  \item Build stage parameters and construction-cost payload for computation.
  \item Execute stage-wise calculations using \texttt{StageCostCalculator}.
  \item Assemble final output:
        \texttt{initial\_stage}, \texttt{use\_stage}, \texttt{reconstruction},
        \texttt{end\_of\_life}.
  \item If requested, generate report with \texttt{core/latex/report.py}.
\end{enumerate}

\subsection{Module Responsibilities}

\renewcommand{\arraystretch}{1.45}
\begin{table}[H]
  \centering
  \caption{Key modules and their responsibilities.}
  \small
  \begin{tabularx}{\linewidth}{%
      >{\columncolor{rowgray}}L{0.40\linewidth}%
      >{\columncolor{rowwhite}}X}
    \toprule
    \rowcolor{navyblue}
    \hc{Module} & \hc{Responsibility} \\
    \midrule
    \texttt{core/main.py}
      & Orchestration, validation call, debug/report mode handling. \\[0.15em]
    \texttt{core/stage\_cost/stage\_cost.py}
      & Full stage-wise economic, environmental, and social computation logic. \\[0.15em]
    \texttt{core/stage\_cost/utils/present\_worth\_factor.py}
      & Present-worth-factor computation for recurring and demolition events. \\[0.15em]
    \texttt{core/latex/report.py}
      & Formula parsing, equation rendering, table rendering, and final
        \texttt{.tex} document generation. \\[0.15em]
    \texttt{inputs/input.py}, \texttt{inputs/input\_global.py}
      & Typed input models and schema-level validation. \\
    \bottomrule
  \end{tabularx}
\end{table}
\renewcommand{\arraystretch}{1}

\screenshotplaceholder
  {Architecture Flow Diagram (Add Screenshot)}
  {Insert diagram: Input $\to$ Main $\to$ StageCostCalculator $\to$ Debug Payloads $\to$ \LaTeX{} Report.}

% =============================================================================
\section{Detailed Cost Logic}

\subsection{Initial Stage}
Initial stage computes:
\begin{itemize}
  \item Initial construction and carbon costs.
  \item Construction-duration road-user and vehicular-emission costs.
  \item Time cost of loan.
\end{itemize}

\begin{equation}
  C_{\text{loan}} = C_{0} \times r \times t_{c} \times I \times \mathrm{SPWF}
  \label{eq:loan}
\end{equation}

where $C_{0}$ is the initial construction cost, $r$ is the interest rate, $t_{c}$ is
the construction duration in years, and $I$ is the investment ratio.

\subsection{Use Stage}
Use stage aggregates routine and major interventions:
\begin{enumerate}
  \item Routine inspection.
  \item Periodic maintenance and periodic carbon impact.
  \item Major inspection.
  \item Major repair (including disruption-linked RUC and vehicular emissions).
  \item Bearing/expansion-joint replacement.
\end{enumerate}

\subsection{Reconstruction Stage}
Applied only when the analysis period exceeds service life. It combines:
\begin{itemize}
  \item Demolition/disposal costs and carbon costs.
  \item Reconstruction costs discounted by reconstruction demolition factor.
  \item Disruption RUC and vehicular emissions for demolition and reconstruction periods.
  \item Time cost of loan for reconstruction and discounted scrap value.
\end{itemize}

\subsection{End-of-Life Stage}
Final demolition factor is used to compute:
\begin{itemize}
  \item Demolition/disposal cost.
  \item Demolition carbon cost.
  \item Demolition-period social and environmental disruption costs.
  \item Final discounted scrap value.
\end{itemize}

\screenshotplaceholder
  {Stage-Wise Output JSON (Add Screenshot)}
  {Insert screenshot of runtime output showing initial/use/reconstruction/end-of-life result blocks.}

% =============================================================================
\section{How the \LaTeX{} Report Is Built}

\subsection{Debug Payload Design}
Each stage breakdown provides three structured data blocks:
\begin{itemize}
  \item \texttt{formulae}
  \item \texttt{inputs}
  \item \texttt{computed\_values}
\end{itemize}
These blocks are consumed by the report engine to generate explainable, traceable
sections.

\subsection{Rendering Strategy}
The report generator:
\begin{enumerate}
  \item Renders title and project parameter tables.
  \item Renders construction-cost inputs.
  \item Renders stage summary tables.
  \item Renders detailed calculations from debug payloads — formula, explanation,
        substitution, result, and abbreviation legend.
  \item Renders cross-stage final summary and optional warnings/notes.
\end{enumerate}

\subsection{Required Section Order (Task Compliance)}

\renewcommand{\arraystretch}{1.35}
\begin{table}[H]
  \centering
  \caption{Mandatory section order in the generated report.}
  \small
  \begin{tabular}{%
      >{\columncolor{rowgray}}C{0.07\linewidth}%
      >{\columncolor{rowwhite}}L{0.52\linewidth}}
    \toprule
    \rowcolor{navyblue}
    \hc{\#} & \hc{Section Title} \\
    \midrule
    1 & Title Page \\
    2 & Construction Cost Inputs \\
    3 & Initial Stage \\
    4 & Use Stage \\
    5 & Reconstruction \\
    6 & End-of-Life Stage \\
    7 & Summary \\
    \bottomrule
  \end{tabular}
\end{table}
\renewcommand{\arraystretch}{1}

\screenshotplaceholder
  {Generated Report Pages (Add Screenshot)}
  {Insert screenshots of one stage summary page and one detailed equation page from the generated PDF.}

% =============================================================================
\section{Requirement-to-Implementation Mapping}

\renewcommand{\arraystretch}{1.45}
\begin{table}[H]
  \centering
  \caption{Traceability from README requirements to implementation locations.}
  \small
  \begin{tabularx}{\linewidth}{%
      >{\columncolor{rowgray}}L{0.42\linewidth}%
      >{\columncolor{rowwhite}}X}
    \toprule
    \rowcolor{navyblue}
    \hc{README Requirement} & \hc{Where It Is Addressed} \\
    \midrule
    Add \texttt{latex\_report} and \texttt{latex\_output\_path} in entry point
      & \texttt{core/main.py::run\_full\_lcc\_analysis()} \\[0.15em]
    Enable internal debug for \LaTeX{} generation without forced debug file dumping
      & \texttt{requested\_debug} vs \texttt{internal\_debug} handling in
        \texttt{main.py} \\[0.15em]
    Generate detailed, human-readable \texttt{.tex} report
      & \texttt{core/latex/report.py::generate\_latex\_report()} \\[0.15em]
    Include formula, inputs, substitutions, step-wise result logic
      & Formula-render and payload-traversal helpers in \texttt{report.py} \\[0.15em]
    Cover all four lifecycle stages and final summary
      & Stage rendering loop + final summary builder in \texttt{report.py} \\[0.15em]
    Add tests
      & \texttt{tests/test\_main.py}, \texttt{tests/test\_latex\_report.py} \\
    \bottomrule
  \end{tabularx}
\end{table}
\renewcommand{\arraystretch}{1}

% =============================================================================
\section{Testing and Evidence}

\subsection{Test Coverage}
Current tests cover:
\begin{itemize}
  \item Debug file behaviour with \texttt{latex\_report=True} and \texttt{debug=False}.
  \item Debug file behaviour with explicit \texttt{debug=True}.
  \item Formula rendering behaviour in report utilities.
  \item Required report section ordering.
\end{itemize}

\subsection{Execution Commands}
\begin{lstlisting}[language=bash]
python -m src.examples.from_dict.example
pytest -q
\end{lstlisting}

\screenshotplaceholder
  {Test Run Output (Add Screenshot)}
  {Insert screenshot of \texttt{pytest -q} output showing all tests passing.}

% =============================================================================
\section{Deliverables Checklist (As Per Task)}

\begin{itemize}[leftmargin=2em]
  \item[$\square$] Updated \texttt{main.py} with \texttt{latex\_report} and
                   \texttt{latex\_output\_path}.
  \item[$\square$] \LaTeX{} report generator module (\texttt{core/latex/report.py}).
  \item[$\square$] At least one test in \texttt{tests/}.
  \item[$\square$] Sample generated \texttt{LCCA\_Report.tex}.
  \item[$\square$] PR with commit history and implementation summary.
  \item[$\square$] Short walkthrough video covering approach, structure, generation
                   flow, and output.
\end{itemize}

\screenshotplaceholder
  {Submission Artifacts (Add Screenshot)}
  {Insert screenshot of repository files, PR, and artifacts to be submitted.}

% =============================================================================
\section{Conclusion}

This implementation directly addresses the README task: it adds report-generation
controls at the main entry point, preserves debug-file behaviour contracts, and
produces an auditable \LaTeX{} report that explains formulas, inputs, substitutions,
and results across all lifecycle stages. The resulting output is suitable for
engineering review and documentation.

\end{document}
